\documentclass[12pt,letterpaper]{article}

% ------------------------------------------------
% PAQUETES
% ------------------------------------------------
\usepackage[utf8]{inputenc}
\usepackage[T1]{fontenc}
\usepackage[spanish,es-nodecimaldot]{babel}
\usepackage{geometry}
\usepackage{array}
\usepackage{longtable}
\usepackage{booktabs}
\usepackage{tabularx}
\usepackage{enumitem}
\usepackage{xcolor}
\usepackage{hyperref}
\usepackage{amsmath,amssymb}
\usepackage{listings}
\usepackage{graphicx}
\usepackage{float}

\geometry{
    left=2.2cm,
    right=2.2cm,
    top=2.2cm,
    bottom=2.2cm
}

% ------------------------------------------------
% CONFIGURACIÓN DE HIPERVÍNCULOS
% ------------------------------------------------
\hypersetup{
    colorlinks=true,
    linkcolor=blue,
    citecolor=blue,
    urlcolor=blue,
    pdftitle={Programa de sesiones: Taller de Herramientas Computacionales},
    pdfauthor={Athziri Hernádnez Jiménez}
}

% ------------------------------------------------
% CONFIGURACIÓN PARA CÓDIGO PYTHON
% ------------------------------------------------
\definecolor{codegray}{RGB}{245,245,245}
\definecolor{codegreen}{RGB}{0,120,80}
\definecolor{codeblue}{RGB}{0,70,150}

\lstset{
    language=Python,
    backgroundcolor=\color{codegray},
    basicstyle=\ttfamily\small,
    keywordstyle=\color{codeblue}\bfseries,
    commentstyle=\color{codegreen},
    stringstyle=\color{red},
    frame=single,
    breaklines=true,
    showstringspaces=false
}

% ------------------------------------------------
% COMANDOS PERSONALIZADOS
% ------------------------------------------------
\newcommand{\python}{\texttt{Python}}
\newcommand{\latexname}{\LaTeX}
\newcommand{\colab}{\textit{Google Colab}}
\newcommand{\vscode}{\textit{Visual Studio Code}}
\newcommand{\github}{\textit{GitHub}}
\newcommand{\overleaf}{\textit{Overleaf}}

\begin{document}

% ------------------------------------------------
% PORTADA
% ------------------------------------------------
\begin{titlepage}
\centering

{\scshape\LARGE Programa de sesiones de clase\par}

\vspace{1cm}

{\Huge\bfseries Taller de Herramientas Computacionales\par}

\vspace{0.5cm}

{\Large Propuesta de organización didáctica para estudiantes de primer semestre de licenciatura de Matemáticas Aplicadas.}
\\[3ex]
{\large \textbf{Athziri Hernández Jiménez}}

\vfill

\begin{tabular}{|ll|} \hline
\textbf{Duración total:} & 64 horas \\
\textbf{Número de sesiones:} & 32 sesiones \\
\textbf{Duración por sesión:} & 2 horas \\
\textbf{Modalidad:} & Teórico-práctica \\
\textbf{Herramientas principales:} &
\python, \colab, \vscode,\\ & \github, \latexname\ y \overleaf\\ \hline 
\end{tabular}

\vfill

\end{titlepage}

\tableofcontents
\newpage

% ==================================================
\section{Presentación del programa}

El presente programa propone una organización progresiva de las sesiones
del curso \textit{Taller de Herramientas Computacionales}, dirigida a
estudiantes de primer semestre de licenciatura en Matemáticas Aplicadas.

La propuesta parte del supuesto de que los estudiantes poseen niveles
heterogéneos de experiencia previa con computadoras y programación.
La secuencia inicia con actividades de familiarización
con el entorno de trabajo y resolución de problemas sencillos, posteriormente avanzar hacia la programación estructurada, el manejo de
datos, la representación gráfica, el cálculo simbólico, aplicaciones
matemáticas y computacionales, a la vez, la elaboración de documentos científicos mediante \latexname. 

Durante el curso se pretenden utilizar cuatro entornos principales:

\begin{itemize}
    \item \textbf{\colab:} para introducir la programación sin requerir
    instalación local y facilitar el trabajo mediante cuadernos interactivos.
    
    \item \textbf{\vscode:} para introducir gradualmente el trabajo con
    archivos de código, scripts, organización de proyectos y ejecución
    local de programas.
    
    \item \textbf{\github:} para introducir principios básicos de control
    de versiones, documentación de proyectos y trabajo reproducible.
    
    \item \textbf{\overleaf:} para la elaboración colaborativa de documentos,
    reportes técnicos, fórmulas matemáticas, tablas, figuras y presentaciones.
\end{itemize}

% ==================================================
\section{Objetivo general}

Introducir al estudiante en el uso informado de herramientas
computacionales para resolver, analizar y comunicar problemas
matemáticos y científicos, mediante el aprendizaje progresivo de
programación en \python, el desarrollo de aplicaciones y experimentos
computacionales, el uso básico de control de versiones con \github\ y
la elaboración de documentos técnicos y científicos mediante
\latexname.

% ==================================================
\section{Objetivos específicos}

Al finalizar el curso, se espera que el estudiante sea capaz de:

\begin{enumerate}[noitemsep]
    \item Reconocer los elementos básicos de un lenguaje de programación.
    
    \item Resolver problemas sencillos mediante algoritmos y programas
    escritos en \python.
    
    \item Utilizar variables, tipos de datos, funciones, estructuras
    condicionales, ciclos, vectores y matrices.
    
    \item Manipular y visualizar datos mediante herramientas
    computacionales.
    
    \item Representar gráficamente funciones y resultados numéricos.
    
    \item Utilizar herramientas básicas de cálculo simbólico.
    
    \item Aplicar métodos computacionales a problemas matemáticos,
    geométricos y científicos.
    
    \item Reconocer la importancia de los errores numéricos, el redondeo
    y la precisión computacional.
    
    \item Generar y utilizar números aleatorios en simulaciones sencillas.
    
    \item Utilizar \github\ como herramienta básica para almacenar,
    organizar y versionar proyectos.
    
    \item Elaborar documentos académicos y científicos utilizando
    \latexname\ en \overleaf.
    
    \item Integrar programación y documentación en un proyecto final.
\end{enumerate}

% ==================================================
\section{Organización general del curso}

El curso se organiza en cuatro bloques principales.
\begin{tabular}{lll}
\toprule
\textbf{Bloque} & \textbf{Temática principal} & \\
\midrule
I & Introducción a la programación y pensamiento computacional &  \\
II & Aplicaciones, visualización y experimentación computacional & \\
III & Introducción a Linux y \github &  \\
IV & Introducción a \latexname\ y comunicación científica & \\
\bottomrule
\end{tabular}


% ==================================================
\section{Estructura sugerida de una sesión de dos horas}

La estructura puede modificarse dependiendo de la complejidad del tema,
pero se propone la siguiente organización general:

\begin{table}[H]
\centering
\begin{tabular}{p{4.5cm}p{2cm}p{9cm}}
\toprule
\textbf{Actividad} & \textbf{Tiempo} & \textbf{Propósito} \\
\midrule
Inicio y recuperación de conocimientos & 10 min &
Relacionar el tema con conocimientos previos y presentar el problema. \\

Introducción conceptual & 20--25 min &
Explicar los conceptos fundamentales mediante ejemplos breves. \\

Demostración guiada & 20--25 min &
Mostrar el procedimiento de programación o uso de la herramienta. \\

Actividad práctica & 35--45 min &
Resolver ejercicios individuales o en parejas. \\

Discusión y revisión & 10--15 min &
Analizar resultados, errores y estrategias de solución. \\

Cierre & 5 min &
Recuperar los conceptos centrales y presentar la actividad siguiente. \\
\bottomrule
\end{tabular}
\end{table}

% ==================================================
\section{Programa de sesiones}
%--------------SESIONES------------------
{\Large{\textbf{Lista de sesiones}}}

\begin{itemize}
    \item Sesión 1. Presentación del curso.
\item Sesión 2. Introducción al pensamiento computacional y primer contacto con Python. Variables y tipos de datos.
 \item Sesión 3.Entrada y salida de información.
 \item Sesión 4. Funciones básicas y organización del código. Estructuras condicionales.
 \item Sesión 5. Ciclos y repetición controlada. 
 \item Sesión 6. Listas y manejo básico de colecciones de datos.
  \item Sesión 7 Práctica integradora previa a la primera evaluación: resolución de problemas utilizando variables, condicionales, ciclos, listas y funciones.
\item Sesión 8. \textbf{ Primera Evaluación: Fundamentos de programación y fundamentos de LaTeX: primer texto.}
%-------------------------------------
\item Sesión 9. Arreglos numéricos con NumPy y operaciones vectorizadas.
 \item Sesión 10. Matrices y arreglos bidimensionales.
 \item Sesión 11. Recursiones
 \item Sesión 12. Cálculo simbólico con SymPy.
 \item Sesión 13. Resolución numérica y gráfica de problemas: raíces y ecuaciones.
 \item Sesión 14. Graficación de funciones. Graficación y visualización de datos y funciones con Matplotlib.
 \item Sesión 15. Errores numéricos, precisión y redondeo.
 \item Sesión 16. Números aleatorios y simulación.
 \item Sesión 17. Representación geométrica y transformaciones en 2D y 3D; Fractales y construcción del Triángulo de Sierpinski. {\textcolor{purple}{Temas comodín}}

\item Sesión 18. Práctica integradora previo a la segunda evaluación: modelación, análisis y visualización de un problema matemático o científico.

\item Sesión 19: \textbf{ Segunda evaluación: aplicación de herramientas computacionales para el análisis y resolución de problemas. LaTeX: modo matemático.}
%-----------------------------------------
 \item Sesión 20. Introducción a Linux y GitHub: repositorios, archivos y organización de proyectos. Control de versiones y organización de proyectos con GitHub.

 %--------------LATEX---------------
 \item Sesión 1. Introducción a LaTeX y Overleaf. Estructura y formato de documentos en LaTeX.
 \item Sesión 2. Listas, tablas y organización de resultados.
 \item Sesión 3. Figuras, gráficas y referencias cruzadas.
 \item Sesión 4. Escritura de expresiones matemáticas en LaTeX.
  \item Sesión 5. Teoremas, Corolarios, demostraciones, etc.
 \item Sesión 6. Referencias bibliográficas y citas.
 \item Sesión 7. Presentación con Beamer.

 
 \item Sesión Final. Proyecto de aplicación computacional.

\end{itemize}

\begin{tabular}{lll} \toprule
\textbf{Calendario}&\textbf{Python} & \textbf{\LaTeX} \\ \midrule
\textbf{Septiembre} &8,22,29 & 3,10,17,24\\ \hline
\textbf{Octubre} & 6,13,15,27 & 1,8, 29\\\hline
\textbf{1era. Evaluación:}& 20 y 22 &\\ \hline 
\textbf{Noviembre} & 3, 5, 10, 12, 17, 19, 24 y 26&\\ \hline
\textbf{Diciembre: 2da. Evaluación:} & 1 y 4 & \\\hline
\textbf{Diciembre: Linux y GitHub}& 8 y 9 &\\ \hline 
\textbf{Enero: Exposiciones proyecto final}& 4, 9, 11 y 16&\\ \hline 
\end{tabular}
\\

% --------------------------------------------------
\subsection{Bloque I. Introducción a Python y pensamiento computacional}

% --------------------------------------------------
\subsubsection*{Sesión 1. Presentación del curso y evaluación de conocimientos generales}

\textbf{Temática:}
\begin{itemize}[noitemsep]
    \item Presentación de la asignatura.
    \item Realización de un cuestionario evaluando los conocimientos generales de los alumnos.   
\end{itemize}

\textbf{Objetivo de la sesión:}

Conocer el estado de conocimiento sobre programación de los alumnos. 


\subsubsection*{Sesión 2a. Introducción al pensamiento computacional y primer contacto con Python.}

\textbf{Temática:}
\begin{itemize}[noitemsep]
    \item Presentación de la asignatura.
    \item ¿Qué es un algoritmo?
    \item ¿Qué es un lenguaje de programación?
    \item Primer contacto con \python\ mediante \colab.
\end{itemize}

\textbf{Objetivo de la sesión:}

Que el estudiante reconozca la relación entre un problema, un algoritmo
y un programa, y ejecute sus primeras instrucciones en \python.\\

\textbf{Ejercicios:}
\begin{enumerate}[noitemsep]
    \item Escribir un algoritmo para preparar una bebida, sandwich, pasta, etc.
    \item Identificar instrucciones ambiguas.
    \item Ejecutar un programa que calcule operaciones aritméticas.
\end{enumerate}

% --------------------------------------------------
\subsubsection*{Sesión 2b. Variables y tipos de datos}

\textbf{Temática:}
\begin{itemize}[noitemsep]
    \item Variables.
    \item Números enteros y reales.
    \item Cadenas de texto.
    \item Operadores aritméticos.
    \item Asignación.
\end{itemize}

\textbf{Ejercicios:}
\begin{enumerate}[noitemsep]
    \item Calcular el área de un rectángulo.
    \item Calcular el perímetro de un círculo.
    \item Convertir una temperatura de grados Celsius a Fahrenheit.
    \item Calcular la distancia recorrida en un movimiento con velocidad
    constante.
\end{enumerate}

\textbf{Actividad práctica integradora:}

Construir un programa que reciba valores definidos por el usuario y
calcule una cantidad física o matemática sencilla.

% --------------------------------------------------
\subsubsection*{Sesión 3. Entrada y salida de información}

\textbf{Temática:}
\begin{itemize}[noitemsep]
    \item Función \texttt{input()}.
    \item Conversión de tipos.
    \item Salida de información.
    \item Formato de resultados.
\end{itemize}

\textbf{Ejercicios:}
\begin{enumerate}[noitemsep]
    \item Calculadora de índice de masa corporal.
    \item Conversor de unidades.
    \item Cálculo del promedio de tres calificaciones.
    \item Cálculo de la ecuación de movimiento rectilíneo uniforme.
\end{enumerate}

% --------------------------------------------------
\subsubsection*{Sesión 4a. Funciones básicas y organización del código}

\textbf{Temática:}
\begin{itemize}[noitemsep]
    \item Funciones matemáticas.
    \item Importación de módulos.
    \item Funciones definidas por el usuario.
    \item Parámetros y valores de retorno.
\end{itemize}

\textbf{Ejercicios:}
\begin{enumerate}[noitemsep]
    \item Crear una función para calcular el área de un triángulo.
    \item Crear una función para convertir kilómetros a metros.
    \item Crear una función para calcular la velocidad media.
    \item Crear una función que determine si un número es positivo,
    negativo o cero.
\end{enumerate}

% --------------------------------------------------
\subsubsection*{Sesión 4b. Estructuras condicionales}

\textbf{Temática:}
\begin{itemize}[noitemsep]
    \item Expresiones booleanas.
    \item Operadores relacionales.
    \item \texttt{if}.
    \item \texttt{if-else}.
    \item \texttt{if-elif-else}.
\end{itemize}
\textbf{Ejercicios:}
\begin{enumerate}[noitemsep]
    \item Determinar si un número es par o impar.
    \item Determinar el mayor de dos números.
    \item Clasificar una temperatura como baja, media o alta.
    \item Resolver un problema de tarifa con diferentes condiciones.
\end{enumerate}

% --------------------------------------------------
\subsubsection*{Sesión 5. Ciclos: repetición controlada}

\textbf{Temática:}
\begin{itemize}[noitemsep]
    \item Concepto de iteración.
    \item Ciclo \texttt{for}.
    \item Función \texttt{range()}.
    \item Ciclo \texttt{while}.
    \item Variables acumuladoras.
    \item Contadores.
    \item Condicionales dentro de ciclos.
    \item Diseño de algoritmos.
\end{itemize}

\textbf{Ejercicios:}
\begin{enumerate}[noitemsep]
    \item Mostrar los números del 1 al 100.
    \item Calcular la suma de los primeros $n$ números.
    \item Generar una tabla de multiplicar.
    \item Aproximar la suma de una serie matemática.
    \item Desarrollar una versión simple de un programa de análisis de calificaciones.
\end{enumerate}

% --------------------------------------------------
\subsubsection*{Sesión 6. Listas y manejo básico de colecciones de datos}

\textbf{Temática:}
\begin{itemize}[noitemsep]
    \item Listas.
    \item Índices.
    \item Acceso y modificación de elementos.
    \item Recorrido de listas.
\end{itemize}

\textbf{Ejercicios:}
\begin{enumerate}[noitemsep]
    \item Calcular el máximo y mínimo de una lista.
    \item Calcular el promedio de una lista.
    \item Contar cuántos valores cumplen una condición.
    \item Construir una lista de cuadrados de números.
\end{enumerate}

% --------------------------------------------------
\subsubsection*{Sesión 7. Práctica Integradora}
%---------------------------------------------------
\subsubsection*{Sesión 8. Primera Evaluación}
%----------------------------------------------------
\subsubsection*{Sesión 9. Arreglos numéricos con NumPy y operaciones vectorizadas}

\textbf{Temática:}
\begin{itemize}[noitemsep]
    \item Introducción a NumPy.
    \item Arreglos unidimensionales.
    \item Operaciones vectorizadas.
    \item Generación de secuencias numéricas.
\end{itemize}

\textbf{Ejercicios:}

Generar valores de posición para:
\begin{enumerate}[noitemsep]
    \item Movimiento rectilíneo uniforme.
    \item Caída libre ideal.
    \item Movimiento con aceleración constante.
\end{enumerate}

% --------------------------------------------------
\subsubsection*{Sesión 10. Matrices y arreglos bidimensionales}

\textbf{Temática:}
\begin{itemize}[noitemsep]
    \item Matrices.
    \item Dimensiones.
    \item Índices.
    \item Operaciones básicas.
    \item Producto matricial.
\end{itemize}

\textbf{Ejercicios:}
\begin{enumerate}[noitemsep]
    \item Crear matrices de diferentes dimensiones.
    \item Calcular la suma de dos matrices.
    \item Calcular el producto de una matriz por un escalar.
    \item Explorar una transformación lineal sencilla.
\end{enumerate}

% --------------------------------------------------
\subsubsection*{Sesión 11a. Funciones, modularidad y reutilización}

\textbf{Temática:}
\begin{itemize}[noitemsep]
    \item Diseño de funciones.
    \item Separación de problemas.
    \item Reutilización de código.
    \item Pruebas básicas.
\end{itemize}

\textbf{Ejercicios:}

Desarrollar una pequeña biblioteca de funciones para:

\begin{itemize}[noitemsep]
    \item Conversión de unidades.
    \item Áreas y perímetros.
    \item Velocidad y aceleración media.
    \item Operaciones estadísticas sencillas.
\end{itemize}
%--------------------------------------------------
\subsubsection*{Sesión 11b. Ceros de una función y raíces de polinomios}

\textbf{Temática:}
\begin{itemize}[noitemsep]
    \item Interpretación gráfica de una raíz.
    \item Raíces de polinomios.
    \item Métodos numéricos sencillos.
    \item Comparación entre resultados simbólicos y numéricos.
\end{itemize}

\textbf{Ejercicios:}
\begin{enumerate}[noitemsep]
    \item Encontrar las raíces de un polinomio cuadrático.
    \item Explorar polinomios de mayor grado.
    \item Graficar la función e identificar sus ceros.
    \item Comparar una solución gráfica, simbólica y numérica.
\end{enumerate}

% --------------------------------------------------
\subsubsection*{Sesión 11c. Recursión}

\textbf{Temática:}
\begin{itemize}[noitemsep]
    \item Concepto de recursión.
    \item Caso base.
    \item Llamada recursiva.
\end{itemize}

\textbf{Ejercicios:}
\begin{enumerate}[noitemsep]
    \item Factorial.
    \item Sucesión de Fibonacci.
    \item Suma de los primeros $n$ enteros.
    \item Comparar una solución iterativa con una recursiva.
\end{enumerate}

% --------------------------------------------------
\subsubsection*{Sesión 12. Cálculo simbólico con SymPy}

\textbf{Temática:}
\begin{itemize}[noitemsep]
    \item Introducción a SymPy.
    \item Variables simbólicas.
    \item Simplificación.
    \item Derivación.
    \item Integración.
    \item Resolución de ecuaciones sencillas.
\end{itemize}



\textbf{Ejercicios:}

Resolver simbólicamente y verificar numéricamente problemas elementales de cálculo.


% ==================================================
\subsection{Bloque II. Aplicaciones y experimentación computacional}

% --------------------------------------------------
\subsubsection*{Sesión 14. Graficación de funciones, Visualización de datos y funciones con Matplotlib}

\textbf{Temática:}
\begin{itemize}[noitemsep]
    \item Introducción a Matplotlib.
    \item Gráficas de funciones.
    \item Ejes, etiquetas y leyendas.
\end{itemize}

\textbf{Actividad práctica:}

Representar gráficamente diferentes funciones y analizar cómo cambian sus gráficas al modificar parámetros.

% --------------------------------------------------
\subsubsection*{Sesión 15. Errores numéricos, precisión y redondeo.}

\textbf{Temática:}
\begin{itemize}[noitemsep]
    \item Precisión finita.
    \item Error absoluto.
    \item Error relativo.
    \item Propagación básica de errores.
    \item Problemas asociados al redondeo.
\end{itemize}

\textbf{Ejercicios:}
\begin{enumerate}[noitemsep]
    \item Comparar aproximaciones de $\pi$.
    \item Analizar el error al redondear valores.
    \item Comparar diferentes representaciones numéricas.
    \item Analizar la diferencia entre dos expresiones matemáticamente
    equivalentes.
\end{enumerate}

% --------------------------------------------------
\subsubsection*{Sesión 16. Números aleatorios y simulación}

\textbf{Temática:}
\begin{itemize}[noitemsep]
    \item Generación de números aleatorios.
    \item Distribuciones básicas.
    \item Simulación computacional.
\end{itemize}

\textbf{Ejercicio:}

Realizar una estimación de una cantidad mediante simulación de Monte Carlo sencilla.
%--------------------------------------------
\subsubsection*{Sesión 17a. Representación geométrica de transformaciones en 2D, {\textcolor{purple}{Optativo}}}

\textbf{Temática:}
\begin{itemize}[noitemsep]
    \item Vectores en el plano.
    \item Transformaciones lineales.
    \item Escalamiento.
    \item Rotación.
    \item Reflexión.
\end{itemize}

\textbf{Ejercicio:}

Representar una figura geométrica y aplicar una matriz de transformación para observar sus cambios.

% --------------------------------------------------
\subsubsection*{Sesión 17b. Transformaciones en 3D. {\textcolor{purple}{Optativo}}}

\textbf{Temática:}
\begin{itemize}[noitemsep]
    \item Coordenadas tridimensionales.
    \item Representación de superficies.
    \item Transformaciones geométricas.
\end{itemize}

\textbf{Ejercicios:}
\begin{enumerate}[noitemsep]
    \item Graficar una superficie sencilla.
    \item Representar un paraboloide.
    \item Explorar una esfera.
    \item Aplicar una transformación a un conjunto de puntos.
\end{enumerate}

% --------------------------------------------------
\subsubsection*{Sesión 17c. Fractales y el Triángulo de Sierpinski, {\textcolor{purple}{Optativo}}}

\textbf{Temática:}
\begin{itemize}[noitemsep]
    \item Autosimilitud.
    \item Algoritmos recursivos.
    \item Construcción computacional de fractales.
\end{itemize}

\textbf{Actividad principal:}

Construcción del Triángulo de Sierpinski mediante una estrategia iterativa o recursiva.

\textbf{Ejercicio de ampliación:}

Explorar otros patrones fractales y comparar el algoritmo utilizado.

% --------------------------------------------------

% --------------------------------------------------
\subsubsection*{Sesión 18: Práctica integradora}

\textbf{Objetivo:}

Integrar los conocimientos de programación, funciones, ciclos, arreglos y visualización.

\textbf{Posibles proyectos:}
\begin{enumerate}[noitemsep]
    \item Simulación de caída libre.
    \item Exploración de trayectorias de un proyectil.
    \item Análisis de una función matemática con parámetros.
    \item Simulación aleatoria.
    \item Generación de un fractal.
    \item Análisis básico de un conjunto de datos.
\end{enumerate}

\textbf{Producto de la sesión:}

Un cuaderno de \colab\ organizado, comentado y funcional.

%------------------------------------------------------
\subsubsection*{Sesión 19: Segunda Evaluación}

% ==================================================
\subsection{Bloque III. Introducción a Linux y GitHub}

%Linux es un sistema operativo de código abierto ampliamente utilizado en entornos científicos, académicos y tecnológicos. Permite al estudiante familiarizarse con la gestión de archivos, la ejecución de programas y el trabajo mediante la terminal o línea de comandos. Una introducción básica a Linux resulta útil para desarrollar autonomía en el uso de herramientas computacionales y comprender mejor los entornos en los que se desarrollan y ejecutan programas científicos.

%La incorporación de \github\ permite introducir a los estudiantes en prácticas básicas de organización y reproducibilidad computacional.
%El objetivo no es convertir al estudiante en un usuario avanzado de control de versiones, sino proporcionar herramientas elementales para guardar, documentar y recuperar diferentes versiones de un proyecto.

% --------------------------------------------------
\subsubsection*{Sesión 20. ¿Qué es Linux y qué es GitHub?}

\textbf{Temática:}
\begin{itemize}[noitemsep]
\item ¿Qué es Linux y qué es un sistema operativo?
\item Distribuciones de Linux y relación con otros sistemas operativos.
\item Estructura básica del sistema de archivos.
\item Introducción a la terminal y la línea de comandos.
\item Navegación entre directorios.
\item Comandos básicos: pwd, ls, cd, mkdir, cp y rm.

    \item Problemas de trabajar con archivos como:
    \texttt{proyecto\_final.py},
    \texttt{proyecto\_final\_ahora\_si.py},
    \texttt{proyecto\_final\_ultimo.py}.
     \item Repositorios.
    \item Archivos y carpetas.
    \item README.
\end{itemize}

\textbf{Actividad práctica Linux:}
\begin{itemize}
\item Exploración del sistema: utilizar los comandos pwd, ls y cd para identificar la ubicación actual, listar los archivos y navegar entre diferentes directorios.
\item Organización de archivos: crear mediante la terminal una carpeta llamada curso\python, y dentro de ella las subcarpetas ejercicios, proyectos y datos.
\item Gestión básica de archivos: crear un archivo de texto, copiarlo a otra carpeta, cambiarle el nombre y finalmente eliminar una copia utilizando comandos de la terminal.
\item Mostrar una situación en la que existen varias versiones de un mismo archivo y discutir las dificultades para identificar cuál contiene la
versión correcta.
\end{itemize}

\textbf{Actividad práctica Github:}
\begin{enumerate}[noitemsep]
    \item Crear una cuenta en \github.
    \item Crear un primer repositorio.
    \item Agregar un archivo \texttt{README.md}.
    \item Describir brevemente un proyecto.
    \item Subir un archivo de \python.
\end{enumerate}

\textbf{Ejercicio:}

Cada estudiante crea un repositorio llamado, por ejemplo, \texttt{taller-herramientas-computacionales} y organiza una estructura
inicial:

\begin{verbatim}
taller-herramientas-computacionales/
|
|-- README.md
|-- python/
|-- ejercicios/
|-- proyecto/
|-- documentos/
\end{verbatim}

% ==================================================
\subsection{Bloque IV. Introducción a LaTeX y comunicación científica}

% --------------------------------------------------
\subsubsection*{Sesión 1. Introducción a LaTeX y Overleaf. Estructura de un documento.}

\textbf{Temática:}
\begin{itemize}[noitemsep]
    \item ¿Qué es \latexname?
    \item Diferencia entre editor visual y sistema de composición. 
    \item Uso inicial de \overleaf. 
    \item Estructura básica de un documento.
    \item Secciones y subsecciones.
    \item Párrafos.
    \item Énfasis tipográfico.
    \item Paquetes.
    \item Comentarios.
    \item Compilación.
\end{itemize}

\textbf{Ejercicio:}

Modificar una plantilla sencilla que contenga:

\begin{itemize}[noitemsep]
    \item Título.
    \item Autor.
    \item Fecha.
    \item Secciones.
    \item Texto.
\end{itemize}
%--------------------------------------------------
\subsubsection*{Sesión 2. Listas, tablas y organización de resultados}

\textbf{Temática:}
\begin{itemize}[noitemsep]
    \item Listas numeradas.
    \item Listas con viñetas.
    \item Tablas.
    \item Alineación de datos.
    \item Uso de tablas para resultados.
\end{itemize}

\textbf{Actividad práctica:}

Construir una tabla de datos generados previamente en \python\ y presentarla correctamente en un documento.

% --------------------------------------------------
\subsubsection*{Sesión 3. Figuras, gráficas y referencias cruzadas}

\textbf{Temática:}
\begin{itemize}[noitemsep]
    \item Inclusión de figuras.
    \item Tamaño y posición.
    \item Leyendas.
    \item Etiquetas.
    \item Referencias cruzadas.
\end{itemize}

\textbf{Actividad integradora:}

Generar una gráfica en \python, guardarla como imagen e incluirla en un
documento de \latexname.

El estudiante deberá explicar:

\begin{enumerate}
    \item Qué representa la gráfica.
    \item Qué variables se presentan.
    \item Qué comportamiento se observa.
\end{enumerate}
% --------------------------------------------------
\subsubsection*{Sesión 4. Escritura de expresiones matemáticas}

\textbf{Temática:}
\begin{itemize} [noitemsep]
    \item Modo matemático.
    \item Superíndices y subíndices.
    \item Fracciones.
    \item Raíces.
    \item Ecuaciones.
    \item Símbolos matemáticos.
\end{itemize}

\textbf{Ejercicios:}

Escribir en \latexname:

\begin{equation}
v = \frac{\Delta x}{\Delta t}
\end{equation}

\begin{equation}
x(t)=x_0+v_0t+\frac{1}{2}at^2
\end{equation}

\begin{equation}
\int_a^b f(x)\,dx
\end{equation}

% --------------------------------------------------
\subsubsection*{Sesión 5. Teoremas, Corolarios, demostraciones, etc.}
% --------------------------------------------------
\subsubsection*{Sesión 6. Referencias bibliográficas y citas}

\textbf{Temática:}
\begin{itemize}[noitemsep]
    \item Importancia de citar fuentes.
    \item Referencias bibliográficas.
    \item Archivos \texttt{.bib}.
    \item Citas dentro del texto.
    \item Introducción a BibTeX y Biber.
\end{itemize}

\textbf{Actividad práctica:}

Crear una pequeña base bibliográfica con tres referencias y citarlas
dentro de un documento.

% --------------------------------------------------
\subsubsection*{Sesión 7. Presentacion Beamer}

\textbf{Temática:}
\begin{itemize}[noitemsep]
    \item Introducción a presentaciones con Beamer.
    \item Estructura de una presentación científica.
    \item Integración de código, resultados y documentación.
\end{itemize}

\textbf{Actividad final:}

Presentar un proyecto breve que integre:

\begin{enumerate}[noitemsep]
    \item Un problema matemático o científico.
    \item Código desarrollado en \python.
    \item Una o más representaciones gráficas.
    \item Un repositorio organizado en \github.
    \item Un reporte escrito en \latexname.
    \item Una presentación breve.
\end{enumerate}

% ==================================================
\section{Propuesta de ejercicios integradores}
{\textcolor{red}{Falta proponer}}
% =------------------------------------------------

\section{Proyecto final}

El proyecto final consiste en realizar una breve investigación sobre un tema de interés relacionado con las matemáticas aplicadas a cualquier área de estudio (ciencias, ingeniería, economía, ciencias sociales, entre otras). El trabajo deberá incluir una propuesta de análisis, modelo o solución en la que se utilicen herramientas de programación —como \textbf{Python}— para el procesamiento, simulación o visualización de resultados. El proyecto deberá documentarse en un \textbf{informe elaborado en LaTeX} y exponerse ante el grupo mediante \textbf{Beamer}.\\

El proyecto deberá permitir identificar claramente:

\begin{itemize}[noitemsep]
    \item Qué problema se estudia.
    \item Qué modelo matemático se utiliza.
    \item Cómo se ejecuta el programa.
    \item Qué resultados se obtienen.
    \item Cómo se interpretan dichos resultados.
    \item Qué fuentes bibliográficas se utilizaron.
\end{itemize}


\textbf{Como mínimo deberá contener:}

\begin{enumerate}[noitemsep]
    \item Planteamiento del problema.
    \item Objetivos del proyecto.
    \item Fundamento teórico
    \item Modelo matemático.
    \item Código comentado.
    \item Resultados numéricos o gráficos.
    \item Análisis de resultados.
    \item  Conclusiones, opiniones finales o estado del arte.
    \item Referencias y citas
    \item Repositorio de \github.
    \item Documento técnico elaborado en \latexname.
    \item Presentación breve en Beamer y oral.
\end{enumerate}

% ==================================================
% ==================================================
\subsection{Producto final esperado y entregables}

Al finalizar el curso, cada estudiante deberá entregar su proyecto computacional organizado de manera similar a:

\begin{verbatim}
proyecto-final/
|
|-- README.md
|
|-- codigo/
|   |-- programa_principal.py
|   |-- funciones.py
|
|-- datos/
|
|-- figuras/
|
|-- documento/
|   |-- reporte.tex
|   |-- referencias.bib
|
|-- presentacion/
    |-- presentacion.tex
\end{verbatim}


Informe con la estructura: 
\begin{itemize}[noitemsep]
    \item Introducción
        \item Fundamentos teóricos 
        \item Metodología 
        \item Resultados
        \item Conclusiones
        \item Bibliografía.
         \item Anexo: Código Python \textbf{(.ipynb y .pdf)} funcional y documentado.
    \end{itemize}
    
El código es de cualquier cosa realizada con python expuesta en su proyecto: Gráficos, resolución de ecuaciones, etc.\\


 \textbf{Presentación Beamer (.pdf y .tex)}\vspace{-2ex}
    \begin{itemize}[noitemsep]
        \item Información sintetizada del informe técnico: Introducción, objetivos, resultados y su análisis, conclusiones y referencias.
        \item Recursos visuales: gráficos, fórmulas, etc. 
        \item Extensión de entre 8 y 12 diapositivas.
    \end{itemize}
    
\textbf{Exposición oral}\vspace{-2ex}

\begin{itemize}[noitemsep]
\item Duración: máximo 10 minutos.
\item Extensión: entre 8 y 12 diapositivas.

Evaluación: se considerará claridad, manejo del tiempo y conocimiento del tema.
\end{itemize}
% ==================================================

\section{Consideraciones didácticas finales}

Dado que se trata de estudiantes de primer semestre, se recomienda
evitar que el curso se convierta únicamente en una enseñanza de sintaxis.
Cada concepto de programación debe asociarse con un problema que el
estudiante pueda comprender, modificar y explorar.

Se propone la siguiente progresión pedagógica:

\begin{enumerate}[noitemsep]
    \item \textbf{Primero:} observar un ejemplo.
    \item \textbf{Después:} modificarlo.
    \item \textbf{Posteriormente:} reconstruirlo con ayuda.
    \item \textbf{Finalmente:} resolver un problema nuevo de manera
    autónoma.
\end{enumerate}

Asimismo, se recomienda mantener una conexión constante entre los
diferentes entornos de trabajo:

\begin{center}
\textbf{Explorar $\rightarrow$ Programar $\rightarrow$ Organizar
$\rightarrow$ Documentar $\rightarrow$ Comunicar}
\end{center}

De esta manera, \colab\ puede utilizarse principalmente durante las
primeras etapas de exploración y aprendizaje; \vscode\ para el
desarrollo progresivo de programas y proyectos; \github\ para la
organización y versionado del trabajo; y \overleaf\ para la elaboración
de reportes y presentaciones.\\[3ex]


% ==================================================
\section{Recursos de consulta recomendados}

Para el desarrollo del curso se recomienda utilizar la documentación
oficial y los recursos educativos correspondientes a:

\begin{itemize}
    \item Python Software Foundation: documentación oficial de Python.
    \item Project Jupyter y Google Colaboratory: documentación sobre
    cuadernos computacionales.
    \item NumPy: documentación oficial para computación numérica.
    \item SymPy: documentación oficial para cálculo simbólico.
    \item Matplotlib: documentación oficial para visualización.
    \item Git: documentación y libro \textit{Pro Git}.
    \item GitHub Docs: documentación sobre repositorios y control de
    versiones.
    \item Overleaf Learn: recursos para aprender \latexname.
    \item Comprehensive TeX Archive Network (CTAN): documentación y
    paquetes de \latexname.
\end{itemize}

% ==================================================


\end{document}